% https://www.amazon.science/_attachment/ara-proposal-template-2025?id=00000195-aa14-ddc9-ad95-fef5abb90000
\documentclass[onecolumn]{article}
\usepackage[a4paper, total={7in, 8in}]{geometry}
\usepackage{hyperref}
\usepackage{amsmath}
\usepackage{amssymb}
\usepackage{pgfgantt}
% \usepackage[authoryear]{natbib}

\usepackage[style=verbose-note,backend=bibtex,maxcitenames=1]{biblatex}
\DeclareNameAlias{sortname}{given-family}
\DeclareNameAlias{default}{given-family}
\renewcommand*{\bibinitdelim}{\addnbthinspace}
\addbibresource{references.bib}

% \bibliographystyle{plainnat}
\usepackage[english]{babel}
% redefine \cite to put citations in footnotes
\let\oldcite\cite
\renewcommand{\cite}[1]{\footnote{\oldcite{#1}}}
\newcommand{\citep}{\parencite}

\hypersetup{
    colorlinks = true, % Enable colored links
    linkcolor  = blue!70!white, % Internal links (e.g., cross-references)
    citecolor  = red!70!white,  % Citations
    urlcolor   = blue!70!white  % External URLs
}

% https://stackoverflow.com/a/4974583/5305365
\usepackage{enumitem}
\setitemize{noitemsep,topsep=0pt,parsep=0pt,partopsep=0pt}


\begin{document}

\newcommand{\bvdecide}{\texttt{bv\_decide}}
\newcommand{\qfbv}{\texttt{QF\_BV}}
\newcommand{\LIA}{\texttt{LIA}}
\newcommand{\Z}{\ensuremath{\mathbb Z}}
\newcommand{\euf}{\ensuremath{\texttt{EUF}}}

\definecolor{almostwhite}{HTML}{efefef}


\title{Toward Certified Machine Learning with Fast, Mechanized Algorithms for Floating Point Verification}
% {Lead: Tobias Grosser, University of Cambridge}}
\date{\vspace{-5ex}}
\maketitle{}
% https://www.amazon.science/_attachment/ara-proposal-template-2025?id=00000195-aa14-ddc9-ad95-fef5abb90000

\section{Principal Investigators}

\begin{itemize}
  \item{PI} Dr. Tobias Grosser, Associate Professor, Computer Laboratory, University of Cambridge
  \item{Co-PI} Siddharth Bhat, PhD Student, Computer Laboratory, University of Cambridge
  \item Cash Funding Needed: £300,000
  \item AWS Promotional Credits needed: £50,000
  \item Amazon contact: Jim Grundy, \texttt{jmgruj@amazon.com}
\end{itemize}

\section{Abstract}


    To mechanize precise, machine-checked definitions of FP8 and
    FP4 in Lean, capturing bit layouts, rounding modes, special values, and
    conversion functions, and implement a new, layered, IEEE784 compliant
    floating-point library. This enables one to develop a formally verified bit-blasting engine in
    Lean that systematically translates FP8/FP4 operations into Boolean
    constraints, proving logical equivalence to high-level semantics, and
    incorporating abstraction/refinement strategies.
    Furthermore, this bitblasting pipeline allows one to mechanize parameterized
    circuit designs for high-performance FP8/FP4
    arithmetic in Lean and prove their behavioral equivalence to the semantic
    model, which yields end-to-end verified hardware modules for the bit-blasting
    pipeline.


\section{Introduction}

\subsubsection{The Critical Need for Trustworthy Hardware Verification} Modern
digital hardware systems, ranging from microprocessors in consumer devices to
specialized accelerators for machine learning and safety-critical control
systems in automotive or aerospace applications, exhibit extraordinary
complexity. Designing and fabricating this hardware is an immensely expensive
undertaking, involving intricate design flows and costly manufacturing
processes. Errors detected late in the design cycle, or worse,
post-fabrication, can lead to catastrophic financial losses, significant delays
in time-to-market, and potential failures in safety-critical scenarios,
jeopardizing lives and infrastructure. Consequently,
rigorous and \emph{trustworthy} verification is not merely beneficial but
absolutely essential. It provides the necessary confidence that a design
behaves according to its specification before committing to the irreversible
and expensive step of manufacturing silicon. High assurance in verification
methodologies is paramount for the economic viability and safe deployment of
advanced hardware.

\subsubsection{The Rise of Custom Low-Precision Arithmetic}
The rapid growth of machine learning and high-performance
computing has spurred the development of specialized hardware accelerators that
utilize custom, low-precision floating-point formats, such as 8-bit (FP8) and
4-bit (FP4)~\cite{micikevicius2022fp8, kim2025investigation}. These formats
deviate from the well-established IEEE-754 standard, offering significant gains
in computational throughput and energy efficiency at the cost of numerical
precision and introducing subtle semantic variations (e.g., in handling
rounding, exponent bias, NaNs, and infinities). Verifying designs incorporating
these novel arithmetic types presents a significant challenge, requiring
precise, formal semantics and reasoning techniques tailored to their specific
behaviours. Effectively applying bit-blasting techniques (a common strategy in
SMT solvers for converting arithmetic problems into boolean satisfiability)
benefits formally verified translations from these
low-precision semantics to Boolean logic.
Our work addresses this emerging need by providing mechanized semantics for FP8/FP4 and developing verified
bit-blasting procedures.

% \subsubsection{Our Contribution: Mechanized and Certified Reasoning for Key Hardware Primitives}
% This project confronts these key challenges by systematically developing fast,
% \emph{mechanized}, and \emph{formally verified} algorithms, data structures,
% and tools within the Lean 4 proof assistant. We aim to build high-assurance
% foundations for finite domain reasoning, directly applicable to hardware model
% checking. Our contributions include verified SMT-based bitvector reasoning
% (WP1), reflexive tactics for core model checking algorithms (WP2), verified
% checkers for industry-standard proof certificates long with novel
% certificate formats (WP3), and a comprehensive framework for verifying
% low-precision floating-point hardware (WP4).

\paragraph{Our Preliminary Work.}

Our preliminary work provides the essential building blocks for the proposed
research, focusing on high-assurance foundations for hardware verification
using formal methods in Lean.

We have made significant contributions to the Lean bitvector library,
implementing the core circuits necessary for bitblasting. This includes adding
support for the SMT-LIB draft standard, including overflow flags, and
implementing efficient overflow detection circuits \cite{gok2001efficient}.
We are developing a formally verified floating point bitblasting
library for Lean, targeting low-precision formats such as 8-bit (FP8) and
4-bit (FP4). These custom formats, used in modern deep learning accelerators,
deviate from the standard IEEE-754 and require precise formal semantics
\cite{micikevicius2022fp8, kim2025investigation}. Our work includes mechanizing these semantics
and developing a verified bit-blasting engine that systematically translates
FP8/FP4 operations into SAT.
Furthermore, we are developing solvers for bitwidth-generalization in Lean.
We develop width-independent decision procedures based on
encoding bitvector problems to finite automata. This exploits connections
between boolean circuits and finite state transducers.
We build proof automation for parametric bitvectors,
forming a cornerstone for our reflexive tactics in model checking.

This body of preliminary work in mechanizing bitvector reasoning,
floating-point semantics, and algorithmic decision procedures within Lean
establishes our expertise and lays a solid groundwork for tackling the
ambitious goals of this project.

\section{Methods}

Recent deep-learning accelerators adopt ultra-low-precision formats such as
8-bit (FP8) and 4-bit (FP4) to trade numerical precision for throughput and
energy efficiency \cite{micikevicius2022fp8}. Unlike standard IEEE-754
types, these custom formats (e.g.\ NVIDIA’s E4M3/E5M2 and the new FP4 variants)
encode exponent bias, mantissa length, and NaN/infinity representations
differently, leading to subtle semantic differences across hardware platforms
\cite{kim2025investigation}. To support sound verification, we will
mechanize precise, machine-checked definitions of FP8 and FP4 in our proof
assistant, capturing bit layouts, rounding modes, special values, and
conversion functions.

Our methodology is to implement a new floating-point library in Lean,
that is IEEE-754 compliant, but with a layered structure enabling opt-in support for NaNs and Infinity.
This will enable us to create ``mixin'' floating point types that provide or omit these features.
Furthermore, each operator will be accompanied by proofs of fundamental properties—such as correct round-to-nearest behavior,
monotonicity, and representable-range bounds—ensuring alignment with the hardware specifications \cite{micikevicius2022fp8}

Bit-blasting reduces high-level arithmetic to Boolean formulas for SAT/SMT
solvers, but correctness of this translation is essential to trust verification
outcomes. While abstraction-driven bit-blasting scales to large bit-widths,
the tiny state spaces of FP8 and FP4 enable exhaustive enumeration of all
bit-patterns. We will develop a formally verified bit-blasting engine in
Lean that systematically translates each FP8/FP4 operation into Boolean
constraints and proves logical equivalence to the high-level semantics.

Concretely, we will specify bit-vector layouts for FP8/FP4 and implement
translation routines that generate SAT/SMT formulas encoding sign, exponent,
and mantissa manipulations. For each operation—such as addition or
multiplication—we will prove a correspondence lemma stating that the Boolean
formula is satisfiable exactly when the intended FP relation holds. We will
also incorporate abstraction/refinement strategies
to prune the search space when full enumeration is unnecessary.

\section{Expected Results}


\section{Funds Needed}

% \newpage
% \appendix
% \section*{Appendix: Project Timeline}
% 
% \begin{ganttchart}[
%     hgrid,
%     vgrid,
%     x unit=0.38cm,
%     y unit title=1cm,
%     y unit chart=1cm,
%     time slot format=simple,
%     group/.append style={fill=almostwhite},
%     group peaks height=0,
%     bar/.append style={fill=black!30!white, draw=none},
%     group/.append style={fill=black!30!white, draw=none},
%     bar label font=\footnotesize\sffamily,
%     group label font=\footnotesize\sffamily,
%     milestone label font=\footnotesize\sffamily,
%     title label font=\footnotesize\sffamily,
%   ]{1}{36}
% 
%    \gantttitle{Project Timeline (36 months)}{36} \\
%     \gantttitlelist{1,...,36}{1} \\
%     \ganttbar{WP1 (Finite Domain Reasoning)}{1}{12} \\
%     \ganttbar{WP2 (Reflexive Model Checking)}{6}{20} \\
%     \ganttbar{WP3 (Certificate Checkers)}{15}{30} \\
%     \ganttbar{WP4 (FP8/FP4 Verification)}{18}{36}
%     % \ganttmilestone{WP1 Done}{8} \\
%     % \ganttmilestone{WP2 Done}{16} \\
%     % \ganttmilestone{WP3 Done}{26} \\
%     % \ganttmilestone{WP4 \& Project Completion}{36}
% \end{ganttchart}
% 
% Our project timeline is carefully designed to ensure proper dependencies and allow for iterative development within the Lean ecosystem. \textbf{WP1} (months 1-12) establishes the foundational bitvector reasoning capabilities, with core functionality delivered by month 8. \textbf{WP2} (months 6-20) begins once WP1's preliminary infrastructure is available, developing reflexive model checking tactics. \textbf{WP3} (months 15-30) leverages the mature foundations from WP1 and WP2 to build verified certificate checkers. \textbf{WP4} (months 18-36) focuses on the specialized domain of low-precision floating-point verification. The staggered approach accommodates the inherent engineering challenges of working with cutting-edge proof assistant technology while ensuring sufficient overlap for knowledge transfer and iterative refinement.
% 
\end{document}
